\section{A uniform upper bound for the extreme slide}
\label{sec:amplitude}

We bound the cohomological degrees in the part of a spherical trace
whose slide eigenvalues have maximal absolute value, as the highest
weight grows along a fixed rational ray. The geometric input is a
uniform upper bound for the costandard tests of
Proposition~\ref{cat:springer}. After semisimple specialization, these
tests control the degrees of fibers of the resonant module, with a
correction determined by compatible nilpotent triples. A finite orbit
decomposition then turns the fiber bounds into the upper estimate of
Theorem~\ref{amp:amplitude}. The module may be neither coherent nor
bounded; the uniformity of the tests is what permits this passage.

We use cohomological grading: $H^i(C[a])=H^{i+a}(C)$.
An upper bound $b$ for a complex means that $H^i=0$ for $i>b$;
it does not assert finite-dimensionality or a lower bound.
Throughout this section all bundle levels, all kernels at other places,
and the number of variable kernels are fixed.  Constants may depend on
these data.  They will not depend on the dominant weights in a nef family.
For a reductive group we either bound the abelian degree or divide by a
fixed lattice of central modifications with cocompact image in the free
abelian degree lattice.  Such a lattice acts freely on that lattice.
The assertions below apply to either convention.

\subsection{Bundle geometry and costandard transforms}

Fix a norm on the semisimple part of the rational cocharacter space.
The instability of a bundle is the norm of its dominant
Harder--Narasimhan type.  The size of a modification is any fixed norm
bound for its relative positions, summed over the geometric legs.
Changing either norm changes only the constants.

\begin{proposition}[Bounded geometry]
\label{amp:geometry}
For a fixed split connected reductive group and fixed finite level,
the following statements hold.
\begin{enumerate}[label=(\alph*)]
\item The number of rational isomorphism classes of bundles of
instability at most $b$ is bounded by $C(1+b)^N$, for fixed $C,N$.
\item There is a locally finite constructible stratification of the
bundle stack such that, on a stratum where the automorphism group has
dimension $r$, the stratum has dimension at most $C-r$.
The same constant works for all Harder--Narasimhan types.
\item Suppose that a sequence of bundles $\mathcal E_n$ has
meromorphic isomorphisms to its point-Frobenius transforms, or
meromorphic automorphisms, with the $n$th modification of size at
most $b_n$.  If the ratio of the instability of $\mathcal E_n$ to
$1+b_n$ is unbounded, a subsequence has a coarsening of the canonical
reduction to one fixed maximal parabolic that these isomorphisms
carry to its point-Frobenius transform, or to itself, respectively.
The assertion holds over geometric points and is
compatible with descent and iteration of the Frobenius isomorphism.
\end{enumerate}
All statements allow an arbitrary effective level divisor, including
nonreduced divisors, and any fixed list of parahoric flags.
\end{proposition}

\begin{proof}
We recall precisely the principal-bundle input.  In arbitrary
characteristic the canonical parabolic reduction exists and is unique;
its Levi bundle is semistable.  Semistable Levi bundles in a fixed
component form a quasicompact stack.  The reduced Harder--Narasimhan
strata are of finite type, and the stack of canonical reductions maps
finitely and radicially onto its stratum.  These are the assertions of
\cite[\S\,2.3.1, Theorem~2.3.3(a),(c), and \S\,2.4.1]{Schieder}.
We use the finite radicial description, not an assertion that this map
is an isomorphism in every characteristic.

For each of the finitely many standard Levis $M$, central cocharacters
cover the free degree lattice of $M$ up to finite index; see
\cite[\S\,2.1.3]{Schieder}.
After tensoring by central bundles, its semistable bundles consequently
belong to finitely many bounded families.  Degree-zero central bundles
are included in these families, since the corresponding Picard stacks
are of finite type.  Fix all the representations of these Levis that
occur in the root filtrations and in highest-weight realizations of
the finitely many maximal flag varieties.  On the bounded families,
the maximal and minimal slopes of these associated vector bundles
are bounded.  Restoring a central translation changes each central
weight quotient by its prescribed scalar degree.  It follows that a
quotient of central weight $\beta$ has all slopes in
\begin{equation}
\label{amp:slope-interval}
 [\langle\beta,\nu\rangle-C_0,
   \langle\beta,\nu\rangle+C_0],
\end{equation}
where $\nu$ is the canonical type and $C_0$ is independent of $\nu$.
There are only finitely many representations here, so one constant
suffices.  This argument uses boundedness of families; it does not
use preservation of semistability by a representation in positive
characteristic.

Choose a fixed large number $A>C_0+2g_X-2$, enlarged for the finitely
many weight filtrations if necessary.  Coarsen the canonical parabolic
by retaining exactly the simple gaps greater than $A$.
Write this parabolic as $P=LU$.  In its induced Levi bundle the
remaining semisimple gaps are bounded by $A$.  It therefore belongs,
modulo central translation in $L$, to a bounded family: its canonical
Levi is one of the semistable Levis considered above, and its remaining
degree gaps range over a finite set of rational lattice points.

Use the relative-root height filtration of $U$, whose successive
quotients are vector groups with linear $L$-action.  This parabolic
filtration and its commutator inclusions hold in every characteristic;
see \cite[Theorem~5.4.3]{ConradReductive}.  Filtering them further
by the weights of the original canonical Levi, every resulting root
has a positive coefficient on a retained gap.  Equation~\eqref{amp:slope-interval}
therefore gives minimal slope greater than $2g_X-2$.
Their $H^1$ vanishes by Serre duality.  The $H^1$ of every successive
vector-group quotient vanishes as well, and the nonabelian cohomology
sequence, applied successively, shows that the $P$-bundle is induced
from its $L$-bundle.  The cutoff was chosen from the original
semistable-Levi families before coarsening, so there is no circular
choice of a bound depending on $A$.

It follows that every bundle is induced from one of finitely many
bounded families for a Levi, followed by a central translation in
that Levi.  The translations required in an HN ball are lattice
points in a ball of radius $O(1+b)$ and are polynomially many.
These translations can be chosen over the fixed finite field after
passing to a fixed finite-index degree lattice: choose a divisor of
positive degree over that field and use its line bundle.  The
remaining degree cosets are absorbed into the bounded families.
The bounded principal-bundle families have finitely many rational
classes.  To see this directly, fix a faithful representation and
a twist for which all the associated vector bundles are globally
generated and have vanishing $H^1$.  A basis of global sections
rigidifies each bundle into one of finitely many finite-type framed
algebraic spaces.  Every rational bundle admits such a basis over
the finite field, and each framed space has finitely many rational
points.  The canonical
reduction of a rational bundle descends by uniqueness, and the
vanishing of $H^1(U)$ above holds over the finite field itself.
Thus every rational class is obtained from the same finite list of
rational Levi families.  This proves (a).

Choose finite-type scheme atlases for these bounded families, with
dimensions bounded by $C$.  Each central translate of each atlas
maps to the bundle stack.  A geometric fiber of such a map over an
object with automorphism dimension $r$ has dimension at least $r$:
it parametrizes isomorphisms from a family member to that object.
The dimension formula bounds the image stratum by $C-r$.
Within each fixed HN stratum only finitely many translating
choices occur, and the HN stratification is locally finite.
Refining the images and the stabilizer dimensions thus proves (b)
with a locally finite stratification.
Universal homeomorphisms in the canonical-reduction description do
not change these dimensions or the constructible-sheaf calculation.

For (c), pass to a subsequence on which a simple gap tends to infinity
relative to $1+b_n$, and fix its maximal parabolic $P_i$.
Choose an actual $G$-equivariant projective embedding
$G/P_i\hookrightarrow\mathbb P(V_i)$ using a sufficiently ample
linearized line bundle with weight a divisible multiple of the
fundamental weight.  In particular the scheme-theoretic stabilizer
of the base-point line is $P_i$; no assertion about the stabilizer
of an arbitrary simple highest-weight module in small characteristic
is required.
The line defining the coarsened canonical reduction is the unique
highest-weight line.  Every other weight differs from this weight
by a positive combination of simple roots involving the selected
root.  Equation~\eqref{amp:slope-interval} shows that the degree of this
line exceeds every slope of the quotient by the selected gap minus
a fixed constant.  A modification of size $b_n$ in this fixed
representation has poles bounded by a divisor of degree $O(b_n)$.
The resulting homomorphism from the highest line to the target
quotient therefore vanishes when the gap is sufficiently large.
The meromorphic isomorphism preserves the reduction line and hence
its maximal-parabolic reduction.  For point Frobenius, the target
canonical type and degrees are the same: point Frobenius acts on
the parameter field, not by Frobenius pullback along $X$.
Uniqueness of the canonical reduction gives the descent and
iteration assertions.

Finally, a full level structure has a space of choices of fixed
finite type and fixed dimension; over the finite field its possible
trivializations form a torsor under the fixed finite group
$G(\mathcal O_D)$.  Parahoric flags have the same boundedness
property.  They change all the preceding constants by fixed amounts.
No argument uses reducedness of $D$.
\end{proof}

\begin{lemma}[The diagonal and the Frobenius graph]
\label{amp:trace-bound}
Let $\mathcal B$ be one of the bundle stacks in
Proposition~\ref{amp:geometry}.  The object
$\Delta_!\mathbf k$ on $\mathcal B\times\mathcal B$ has a uniform
upper perverse bound.  Moreover, for a constructible complex
$\mathcal F$ on $\mathcal B\times\mathcal B$ with upper perverse bound
$N$, pullback to the point-Frobenius graph followed by compact
cohomology has upper bound $N+C_1$, where $C_1$ is independent of
$\mathcal F$.
\end{lemma}

\begin{proof}
We use middle perversity with smooth descent, as in
\cite[Lemma~4.1 and Theorems~4.2, 5.1]{LaszloOlssonPerverse}, so that
$\mathbf k[-r]$ is perverse on the classifying stack of a smooth
group of dimension $r$.  Stratify further by stabilizer dimension
and by the isomorphism relation.  Over a pair of isomorphic bundles
with stabilizer dimension $r$, the stalk of $\Delta_!\mathbf k$
is the compact cohomology of their isomorphism scheme and has upper
degree $2r$.  The corresponding locus of isomorphic pairs has
dimension at most $C-2r$, by Proposition~\ref{amp:geometry}(b).
The support criterion for the upper perverse truncation gives a
bound independent of $r$.

At a point of the Frobenius graph the stabilizer dimension in
$\mathcal B\times\mathcal B$ is $2r$.  Every stratum through it has
dimension at least $-2r$.  Thus the stalk upper bound of
$\mathcal F$ is at most $N+2r$.  On the graph the relevant stratum
has dimension at most $C-r$.  Compact cohomology of a sheaf on a
finite-type algebraic stack of dimension $a$ vanishes in degrees
above $2a$; see \cite[Theorem~1.4]{SunStacks}.
The resulting bound is
\[
 (N+2r)+2(C-r)=N+2C.
\]
Stratification and the spectral sequence for cohomology sheaves give
the same estimate for complexes.  Apply the calculation on
finite-type opens and then take the defining filtered colimit.
The bounds are independent of the open, so they persist in that
colimit.  The central-degree convention is treated componentwise.
\end{proof}

\begin{proposition}[Uniform geometric tests]
\label{amp:tests}
In the hyperspecial cyclic calculation, a fixed list of perfect
spherical labels has bounded-above cohomology.  The same upper bound
property holds uniformly when the list also contains a fixed number
of Springer pushforwards
\[
 R\pi_*\mathcal O_{\widetilde{\mathcal N}}(\eta_j),
 \qquad
 \pi:\widetilde{\mathcal N}\longrightarrow\gd,
\]
whose line bundles are independently nef.  Fixed additional twists
or fixed perfect factors are allowed.  The number of factors and
the additional twists are fixed; their nef weights may vary
independently without changing the upper bound.
\end{proposition}

\begin{proof}
Compute the trace by applying the Hecke transforms to one factor
of $\Delta_!\mathbf k$ and then applying
Lemma~\ref{amp:trace-bound}.  Every fixed bounded kernel has finite
upper perverse amplitude.  A costandard kernel has a bound
independent of its length, as follows.

Let $\mathcal H_w$ be its exact-position Hecke stratum, with source
and target maps $p,q$.  The map $p$ is smooth of relative dimension
$\ell(w)$ and $q$ is affine; these facts are checked after smooth
local trivialization, where their fibers are the affine Schubert
cell.  In the closure correspondence the target map is proper.
Here is the coefficient comparison for an arbitrary locally bounded
constructible input.  Choose a sufficiently deep jet subgroup acting
trivially on the chosen Schubert closure, and pass to its smooth
torsor of trivializations.  The open and closed Hecke spaces then
become products over the source.  For bounded constructible
complexes the formula
\[
 R(\operatorname{id}\times j)_*(A\boxtimes B)
 \simeq A\boxtimes Rj_*B
\]
follows by Verdier duality from the corresponding extension-by-zero
formula.  It does not require $A$ to be lisse.  Thus smooth base
change \cite[\S\,12.1]{LaszloOlssonSixII} and descent identify
the costandard coefficient with
$Rj_*p^*(-)[\ell(w)]$.  Proper pushforward from the closure
identifies the costandard transform with
\[
 Rq_*p^*(-)[\ell(w)].
\]
Smooth pullback with this shift preserves perversity, and affine
Artin vanishing says that $Rq_*$ preserves an upper perverse bound;
see \cite[\S\,3.2, paragraph preceding Theorem~3.2]{BeilinsonPerverse}.
Both assertions descend from the smooth scheme charts, since the
maps of the correspondence are representable.  The resulting upper
bound is independent of $w$.

Passage between the fixed parahoric levels costs only a fixed
amplitude, because their fibers are fixed flag varieties.
Forgetting and then right-averaging torus equivariance has a
representable torus-torsor fiber.  Its direct image has a finite
Postnikov filtration with constant cohomology pieces $H^i(T,\mathbf k)$:
translations act trivially on torus cohomology.  This is cohomology
of $T$, not of $BT$, and no splitting of the filtration is needed.
It adds only fixed shifts.  This is the averaging convention in
Proposition~\ref{cat:springer}.
By the same proposition, nef Springer line bundles correspond to
costandard Wakimoto translations, with shifts independent of their
weights.  A bounded additional translation can be composed with a
nef translation at a fixed cost.  Convolution of a fixed number of
these transforms preserves the uniform upper bound just proved.
The spherical polynomial action and the identification of their
products with derived tensor products are those of
Theorem~\ref{cat:satake} and Proposition~\ref{cat:springer}.
\end{proof}

\subsection{Specialization and uniform Levi tests}

Fix a closed leg of degree $d$, set $Q=q^d$, and choose a complex
realization of $\mathbf k$.  Let $t\in\Gd$ be semisimple and put
\[
 H=Z_{\Gd}(t),\qquad
 E_t=\{e\in\gd:\operatorname{Ad}(t)e=Qe\}.
\]
In a torus containing $t$, decompose $\log_Q|t|$ into its
connected-central and semisimple components, using the rational
direct sum of the connected-center and derived-torus cocharacter
spaces.  Throughout this section $\lambda$ denotes the
\emph{semisimple component}, and we assume that $\lambda$ is
rational.  No rationality is required of the connected-central
component.  This convention applies to every grading and
cocharacter below, as well as to metric pairings.
The direction $\lambda$ is central in $H$: conjugation fixing $t$
also fixes its absolute-value direction and its semisimple
projection.  Roots annihilate the discarded central component;
therefore every vector in $E_t$ still has $\lambda$-weight one.
Replacing $t$ by $tz$, for $z\in Z(\Gd)^\circ$, leaves
$\lambda$ unchanged.

Specialization at the class of $t$ always means the following finite
operation. Choose once algebra generators $a_1,\ldots,a_r$ of
$\mathbf k[\Gd]^{\Gd}$, and tensor the cyclic calculation with
the Koszul complex of $a_1-a_1(t),\ldots,a_r-a_r(t)$. These functions
generate the maximal ideal of $[t]$. No regularity of the adjoint
quotient is assumed. The generators have degree zero, and the number
and degrees of the Koszul factors are independent of $t$. We retain
all the derived structure produced by this operation.

\begin{lemma}[The specialized test module]
\label{amp:specialization}
After this specialization, the cyclic calculation is a pairing
with an $H$-equivariant dg module $M$ over
$\operatorname{Sym}(E_t^*[-2])$.  The module need not be coherent
or bounded.  The action of a spherical polynomial at any one
geometric factor is restriction of its linear functions to $E_t$.
The uniform tests of Proposition~\ref{amp:tests} remain valid.
They also allow fixed equivariant perfect factors on the flag
varieties, in particular restriction to fixed closed
$H^\circ$-orbits of Borels.
\end{lemma}

\begin{proof}
By Proposition~\ref{cat:loops}, the holonomy satisfies the derived
linear equation $\operatorname{Ad}(s)e=Qe$.
The fiber of the adjoint quotient over $[t]$ consists of elements
conjugate to $tu$, with $u$ unipotent in $H$.
Indeed the adjoint quotient is the Weyl quotient of a maximal torus
\cite[Theorem~4.1\textup{(2)}]{LeeAdjointQuotients}, and representation
characters are unchanged by the commuting unipotent factor in Jordan
decomposition.
Here is an algebraic slice description, with no completion of the
coefficient module.  Let $W$ be the sum of the nonresonant
$\operatorname{Ad}(t)$-eigenspaces in $\gd$.  Choose an
$H$-stable affine open $U\subset H^\circ$ containing the unipotent
locus on which $1-\operatorname{Ad}(tu)$ on
$\gd/\operatorname{Lie}H$ and $\operatorname{Ad}(tu)-Q$ on
$W$ are invertible.  Both determinants are nonzero at every
unipotent $u$.  There are only finitely many $H$-classes of
semisimple $u$ for which $tu$ is conjugate to $t$, since their
torus representatives come from the finite Weyl orbit of $t$.
Invariant functions allow us also to remove all these classes
except $u=1$.
The conjugation map
\[
 \Gd\mathbin{\times}^{H}(tU)\longrightarrow\Gd
\]
is etale: its transverse differential is the first displayed
invertible operator.  After the removal of the other sheets,
uniqueness of the Jordan semisimple part identifies every
conjugating ambiguity over the class fiber with $H$.
The restriction to the schematic class fiber is therefore etale
and universally bijective, hence an isomorphism.
Being etale, it is an isomorphism over every infinitesimal
thickening of that fiber as well.  Flat pullback and descent
therefore identify the supported quasi-coherent complexes on
these neighborhoods, equivariantly.  This assertion applies to
arbitrary complexes: on affine charts it holds for modules killed
by each power of the fiber ideal, hence for their filtered unions;
flatness computes pullback on cohomology in every degree.
The finite Koszul specialization has cohomology annihilated by
the fiber ideal, so it is covered by this argument.  In particular,
all equivariance remains ordinary rational $H$-equivariance.

Over $\mathcal O(U)$, the invertible nonresonant block of
$\operatorname{Ad}(tu)-Q$ gives a change of odd Koszul generators
for which their differentials are precisely the coordinates of
$W$.  Eliminating these contractible pairs is an actual dg-algebra
quasi-isomorphism over $\operatorname{Sym}(E_t^*[-2])$.
Solve successively for the transported variables at the other
geometric factors in the same way.  Push forward the unipotent
and additional derived variables along their affine projection.
They remain in $M$, including all invariant Koszul generators and
the residual equation $\operatorname{Ad}(u)e=e$.
The polynomial-action assertion follows by solving the equations
starting at the chosen geometric factor.  The description is
$H$-equivariant.

Tensoring with the finite Koszul complex changes upper bounds by
only a fixed amount.  For a fixed finite-dimensional representation
of a reductive subgroup of $\Gd$, every irreducible summand occurs
in a restriction from $\Gd$.  One proof uses the surjection
$\mathbf k[\Gd]\to\mathbf k[H]$, takes a finite-dimensional
$\Gd$-subrepresentation containing lifts of its matrix coefficients,
and then uses complete reducibility.  Thus all fixed $H$-tests are
available.  Induction from $H^\circ$ to $H$ is finite, so it also
allows $H^\circ$-invariants and fixed $H^\circ$-tests.
These retracts are taken after specialization on the slice; their
maps need not extend equivariantly before specialization.

Here is the perfect-generation detail for the flag factors.
On the smooth projective flag variety choose an ample equivariant
line bundle $A$.  Equivariant Serre evaluation gives surjections
from finite sums of $V\otimes A^{-n}$ onto any equivariant coherent
sheaf.  For an explicit finiteness argument, let the flag variety have
dimension $a$.  More than $2a$ successive evaluations give an
extension whose class vanishes, since the local-to-global Ext
spectral sequence and exact reductive invariants give
$\operatorname{Ext}^{n}_{H^\circ}(F,K)=0$ for coherent $F,K$ and
$n>2a$.  The truncated evaluation complex therefore contains $F$
as a retract.  Applying finite cohomology filtrations shows that
every perfect complex is in the finite thick closure of these
bundles.  Apply this statement to its dual and dualize: the same
perfect complex is in the finite thick closure of bundles
$V\otimes A^n$, $n\geq0$.  All representations, twists, cones,
and shifts in this construction are fixed once the perfect complex
is fixed.  Tensoring a variable nef line with these positive twists
keeps it nef.  The structure sheaf of a fixed closed smooth
$H^\circ$-flag orbit is perfect on the ambient flag variety, so the
assertion applies to its restriction test.  This proves the last
claim with constants independent of the variable nef weights.
\end{proof}

The specialized module $M$ satisfies the geometric tests, but need
not be coherent or bounded. Bounds for each fixed representation
test alone would not control cohomology
whose representation types vary with the degree. We next obtain
bounds uniform under character twists on each Levi. In
Proposition~\ref{amp:fiber}, these twists and a finite list of fixed
representations will detect every representation constituent, turning
the invariant test bounds into a bound on the whole fiber.

For a cocharacter $\gamma$ commuting with $t$, define
\[
 L=Z_{\Gd}(\gamma),\qquad H_L=H^\circ\cap L,
 \qquad E_L=E_t\cap\operatorname{Lie}L,
 \qquad M_L=M\otimes^{\mathbf L}_{\mathbf k[E_t]}\mathbf k[E_L].
\]
The group $H_L$ is connected reductive, by the structure of torus
centralizers \cite[Theorem~17.38 and Corollary~17.59]{MilneAlgebraicGroups}.
All coordinate rings in
this notation carry the degree-two shearing.  A representation
test is placed in ordinary cohomological degree zero.
Since $\lambda$ is central in $H$, it commutes with $\gamma$
and is a rational central cocharacter of $H_L$; its weight
on $E_L$ is one.

\begin{proposition}[Uniform central-character tests]
\label{amp:levi}
For every such $L$ and every fixed finite-dimensional rational
$H_L$-representation $U$, there is a constant $C(U,L)$ such that
\begin{equation}
\label{amp:levi-bound}
 H^i\bigl((M_L\otimes U\otimes\eta)^{H_L}\bigr)=0
 \quad\text{for }i>C(U,L),\qquad \eta\in X^*(L).
\end{equation}
Invariants denote invariant global sections.  The constant is
independent of the character $\eta$.
\end{proposition}

\begin{proof}
Choose Borels containing a reference torus with $t,\gamma$.
First order their roots by $\gamma$ and $-\gamma$, respectively;
within $L$, order them positively on $\lambda$ and complete the
order arbitrarily on the zero-weight roots.  Their
$H^\circ$-orbits are complete flag varieties.  Use two Springer
tests restricted to these orbits, as permitted by
Lemma~\ref{amp:specialization}.

Let $P_\gamma,P_{-\gamma}$ be the two opposite parabolics.
Choose an integral character $\chi$ strictly dominant for the first
partial flag variety and zero on the coroots of $L$.
The extremal pairing in mutually dual representations gives an
invariant section of the product of two nef line bundles on the
partial flags. Project the selected complete-flag orbits to these
partial-flag varieties. On the product of the two resulting
$H^\circ$-orbits, the section's nonvanishing locus is the locus
of opposite partial flags, which is
$H^\circ/H_L$. Indeed the ambient big-cell Gauss factorization
is unique and is preserved by conjugation by $t$.  If an element
of $H^\circ$ admits that factorization, its unipotent and Levi
factors centralize $t$ as well and lie in the corresponding
parabolics of $H^\circ$.  Thus ambient oppositeness restricts to
exactly the opposition locus for $H^\circ$.  Localizing by that section is the filtered colimit
of simultaneous positive twists by the two nef lines.
The section and its localization maps have degree zero.

The pairs of complete flags form a fibration over this opposition
locus. Over the reference opposite partial-flag pair, its fiber is
the product of two complete flag varieties for $H_L$. The resonant
parts of the two ambient nilradicals are
\[
 E_{\gamma\geq0}\quad\text{and}\quad E_{\gamma\leq0}.
\]
Indeed the zero $\gamma$-part $E_L$ has positive $\lambda$-weight
and belongs to both nilradicals, independently of the remaining
flags; the roots of $H_L$ have $\lambda$-weight zero.
The displayed subspaces span $E_t$ and intersect in $E_L$, so
\[
 \mathbf k[E_{\gamma\geq0}]
   \otimes^{\mathbf L}_{\mathbf k[E_t]}
 \mathbf k[E_{\gamma\leq0}]
 \simeq \mathbf k[E_L].
\]
Thus their derived intersection over $E_t$ has no excess.

For completeness, the other torus eigenspaces do not obstruct
splitting off this intersection.  Decompose the ambient Lie
algebra as $E_t\oplus W$ by $t$-weights.  The projectors are
polynomials in $\operatorname{Ad}(t)$, so this also splits the
universal nilradical subbundles globally and equivariantly.  Its Koszul base change to $E_t$ has zero differential
on the factors from $W$, so these factors are finite exterior
algebras of equivariant vector bundles.  Their exterior-degree-zero
summand is canonical; no splitting of the varying nonresonant
subbundle inside $W$ is required.  Taking this summand in both tests leaves
exactly $\mathbf k[E_L]$ on the opposite-pair slice. Tensoring with
$M$ therefore gives $M_L$. The remaining
complete flags contribute
$R\Gamma(\mathcal O)=\mathbf k$.  By equivariant descent from
$H^\circ/H_L$, the resulting invariant test is the expression in
Equation~\eqref{amp:levi-bound}.

We check the uniformity in $\eta$.  Give the first partial-flag line
fiber character $\eta$.  After the $n$th localization twist the two
fiber characters are, with the opposite-Borel conventions,
\[
 \eta-n\chi,\qquad n\chi.
\]
Both lines are nef for every sufficiently large $n$, since $\chi$
is strictly positive on precisely the omitted simple coroots and
$\eta$ vanishes on the Levi coroots.  The starting value of $n$ may
depend on $\eta$.  It changes neither the localization colimit nor
the bound: every term in this cofinal tail is in the same uniformly
bounded nef family of Proposition~\ref{amp:tests}.
The fixed resolutions in Lemma~\ref{amp:specialization} introduce
only fixed additional positive twists and fixed shifts.
The character $\eta$ and the localization characters are trivial
on the derived subgroup of $L$, so the residual complete flags still
push their structure sheaves to constants.
If $L=\Gd$, the character line is underlying-trivial and already
nef; this also covers the case with no partial flags.

A fixed $U$ is obtained from a fixed representation test by the
restriction and summand argument in
Lemma~\ref{amp:specialization}.  Finally, global sections on these
fixed flag varieties have finite cohomological dimension,
reductive invariants are exact, and filtered colimits of vector
spaces are exact.  All the operations and the direct-summand
projection above are therefore valid for arbitrary dg $M$.
They preserve a uniform upper bound, proving the proposition.
\end{proof}

\subsection{Orbits, corrected fibers, and local cohomology}

For $e\in E_L$, choose a Jacobson--Morozov triple $(e,h_e,f)$
inside $\operatorname{Lie}L$ such that
\[
 \operatorname{Ad}(t)(e,h_e,f)=(Qe,h_e,Q^{-1}f).
\]
We use $h_e$ for the integral neutral cocharacter as well as its
differential.  Conjugate it into the reference torus of $H_L$ when
computing pairings.

\begin{lemma}[Resonance orbits]
\label{amp:orbits}
Every element of $E_L$ is nilpotent and admits such a compatible
triple.  There are finitely many $H_L$-orbits in $E_L$.
For $e\in E_L$, every $h_e$-weight in the normal space
\[
 N_e=E_L/[e,\operatorname{Lie}H_L]
\]
is nonpositive.
\end{lemma}

\begin{proof}
Choose a faithful algebraic representation
$\rho:L\hookrightarrow\operatorname{GL}(V)$.
The relation $\operatorname{Ad}(t)e=Qe$ says that
$d\rho(e)$ is conjugate to $Q\,d\rho(e)$.
Its finite multiset of eigenvalues is therefore preserved by
multiplication by $Q$. Since $Q>1$, every eigenvalue is zero.
Thus $d\rho(e)$, and hence $e$, is nilpotent.
Start with the Jacobson--Morozov
theorem in characteristic zero
\cite[Theorems~3.1 and~3.2]{BMYJM}.
On the torsor of triples with first member $e$, let $(t,Q)$ act by
\[
 (e,h,f)\longmapsto
 (Q^{-1}\operatorname{Ad}(t)e,\operatorname{Ad}(t)h,
                  Q\operatorname{Ad}(t)f).
\]
The closure of the subgroup generated by $(t,Q)$ is diagonalizable.
The torsor is under the unipotent radical of the centralizer of $e$
\cite[Proposition~3.3]{BMYJM};
successive vector-group quotients have vanishing first cohomology
for a diagonalizable group in characteristic zero.  It therefore
has a fixed point, which is the required compatible triple.

Put $c=tQ^{-h_e/2}$.  It commutes with the triple.
Decompose $\operatorname{Lie}L$ into simultaneous $c$-eigenspaces
and irreducible $\mathfrak{sl}_2$-modules $V_n$.
A vector of $h_e$-weight $k$ belongs to $E_L$ precisely when its
$c$-eigenvalue is $Q^{1-k/2}$.
Its predecessor of weight $k-2$, when present, has $t$-eigenvalue
one and lies in $\operatorname{Lie}H_L$.
The raising map is onto this vector unless $k=-n$ is the lowest
weight.  Hence every normal weight is of the form $-n\leq0$.

For finiteness, conjugate $h_e$ into the fixed torus by $H_L$.
Its central projection in $L$ is zero, because the triple comes
from $\mathrm{SL}_2$.  Its root pairings are integral and bounded
in absolute value by $\dim\operatorname{Lie}L$: they are weights
of an $\mathfrak{sl}_2$-representation of that dimension.
There are consequently only finitely many possibilities in the
fixed cocharacter lattice.  Fix one such $h$.  On the $c=1$ part of the same calculation the
raising map gives a surjection
\[
 \operatorname{Lie}Z_{H_L}(h)
 \longrightarrow E_L\cap(\operatorname{Lie}L)_2(h).
\]
Thus every $e$ admitting this $h$ has an open
$Z_{H_L}(h)$-orbit in that vector space.  An irreducible vector
space has at most one open orbit.  The finite list of $h$'s
therefore gives finitely many $H_L$-orbits.
\end{proof}

We make the grading used for a nonzero fiber explicit.
Let $a=2\lambda$ and first deshear by adding the $a$-weight to
cohomological degree.  A linear function on $E_L$ has original
degree $2$ and $a$-weight $-2$, so the coordinate ring now has
degree zero.  Take the derived fiber at $e$ in this grading.
The rational cocharacter
\[
 \rho_e=h_e-a
\]
fixes $e$, and hence belongs to its stabilizer $S_e=(H_L)_e$.
Add its weight to the fiber degree.  The result is called the
\emph{corrected fiber}, denoted $F_e^{\mathrm{corr}}$.
Equivalently, the total correction is by $h_e$.
Here is a literal implementation for rational $\lambda$.
Choose $N$ for which $N\lambda$ is integral, and decompose into the
finitely many sectors of $\lambda$-weights $w$ modulo $\mathbb Z$.
In the sector represented by $r\in[0,1)\cap N^{-1}\mathbb Z$,
replace degree $i$ by the integral degree $i+2(w-r)$ and retain
the numerical offset $2r$.  The actual shifts are even integers,
so all dg signs are unchanged.  A linear function changes
$(i,w)$ to $(i+2,w-1)$ and has new degree zero.
Tensor products add offsets, with an even integral carry.
Thus deshearing is a finite collection of honest complexes with
rational numerical offsets. Write $W_{e,r}$ for the ordinary derived
fiber in the sector represented by $r$. We use $H^j(W_e)$ for
sectorwise cohomology in numerical degree $j$: a class in
$H^n(W_{e,r})$ has $j=n+2r$. The actual complexes retain integral
degrees. The assertion that the corrected fiber has upper bound $B$
means precisely
\begin{equation}
\label{amp:corrected-fiber}
 H^j(W_e)_b=0\quad\text{whenever }j+b>B,
 \qquad b=\text{$\rho_e$-weight}.
\end{equation}
All nonzero-point pullbacks are taken after this deshearing.
Differentials retain degree one and preserve $b$.
Apply the same deshearing to every representation test
$U\otimes\eta$.  On each of its simultaneous weight spaces,
the desheared degree plus the $\rho_e$-weight is the
$h_e$-weight, since $a+\rho_e=h_e$.
In particular, the possible $a$-weight of a character
$\eta\in X^*(L)$ cancels its $\rho_e$-weight: its $h_e$-weight
is zero because the triple lies in the derived group of $L$.
On $H_L$-invariants the total $a$-weight is zero, so deshearing
does not change the output cohomological degrees.

\begin{lemma}[Contributions of an orbit]
\label{amp:support}
Suppose $F_e^{\mathrm{corr}}$ has upper bound $B_e$.
For every finite-dimensional $H_L$-representation $V$, the
contribution of the orbit $H_Le$ to invariant sections of
$M_L\otimes V$ has upper bound
\begin{equation}
\label{amp:orbit-bound}
 B_e+C_e+\max\{\text{$h_e$-weights of }V\},
\end{equation}
where $C_e$ is independent of $V$.
This assertion uses local cohomology in an open subset where the
orbit is closed and is valid for arbitrary dg $M_L$.
\end{lemma}

\begin{proof}
Deshear as above.  In an invariant open subset where the orbit is
closed it is a smooth regular immersion into the smooth space
$E_L$.  If $\mathcal I$ is its ideal, local cohomology is
\[
 R\Gamma_{H_Le}(M_L)
 =\operatorname*{colim}_{p}
 R\mathcal Hom(\mathcal O/\mathcal I^p,M_L).
\]
Each quotient is perfect on the regular ambient space.
Its finite filtration has graded pieces the symmetric powers of
the conormal bundle.  Regular-immersion duality therefore expresses
the corresponding support terms as
\[
 i_*\bigl(Li^*M_L\otimes\det N\otimes
                      \operatorname{Sym}^{a}N\bigr)[-c],
 \qquad c=\operatorname{codim}(H_Le),
\]
where $M_L$ is read in the desheared grading and $i$ denotes the
regular immersion.  The shift $[-c]$ raises degrees by $c$;
the normal symmetric powers, rather than the conormal ones, are
the factors relevant to the upper bound.

The action of $\rho_e$ on a normal vector of $h_e$-weight $k$
has weight $k-2$.  By Lemma~\ref{amp:orbits} this is at most $-2$.
Consequently arbitrary normal powers cannot increase the corrected
upper bound.  The determinant and the codimension shift cost a
constant depending only on the orbit.  This establishes a bound
uniform in $p$ before passage to the filtered colimit.

Equivariant sections on the orbit are rational group cohomology
of $S_e$.  Its unipotent radical is computed by the finite exterior
Lie-algebra cochain complex, and invariants for a reductive Levi
are exact.  The triple centralizer is a Levi in the centralizer
of $e$ \cite[Proposition~3.3]{BMYJM}; taking the fixed points of
the compatible diagonalizable action gives the same decomposition
for $S_e$.  The cocharacter $\rho_e$ lies in, and is central in,
this triple-centralizer Levi.  If $\mathfrak u$ is the Lie algebra
of the unipotent radical, a possible additional cost is
\[
 \max_{0\leq a\leq\dim\mathfrak u}
 \left(a+\max\operatorname{wt}_{\rho_e}
                          \bigwedge^a\mathfrak u^*\right).
\]
The exterior complex has finitely many columns, so it is valid
for unbounded coefficients without an infinite-totalization issue.
Finite component invariants are exact in characteristic zero.
The test $V$ is desheared too: its $a$-weight-$\alpha$ space
is placed in numerical degree $\alpha$.  Its numerical degree
plus its $\rho_e$-weight is therefore its $h_e$-weight.
Taking invariants of the reductive triple-centralizer Levi forces
total $\rho_e$-weight zero.  On the original $H_L$-invariants
the total $a$-weight is zero as well, so the resulting numerical
degree is the original cohomological degree.
This gives Equation~\eqref{amp:orbit-bound}.

No finite-generation hypothesis has been used.  Tensoring with the
perfect quotients and their finite filtrations is legitimate for
unbounded dg modules.  The operation displayed above is a filtered
colimit, not an inverse limit.  Exactness of filtered colimits
preserves the upper bound uniform in $p$.  Finally there are
finitely many orbits, so their support triangles give a finite
devissage.
\end{proof}

\begin{proposition}[Bounded corrected fibers]
\label{amp:fiber}
For every allowed Levi $L$ and every $e\in E_L$, the corrected
fiber $F_e^{\mathrm{corr}}$ is bounded above.
The bound is a bound on the whole complex; its cohomology need
not be finite-dimensional.
\end{proposition}

\begin{proof}
Induct on the semisimple rank of $L$, using
Proposition~\ref{amp:levi} for all its smaller allowed Levis.
If the centralizer of the triple in $H_L$ has a torus noncentral
in $L$, choose a cocharacter in that torus noncentral in $L$.
Its centralizer is a proper Levi $L'$ of $L$ containing the triple.
It is still allowed: combine this cocharacter with $\gamma$,
using a sufficiently large multiple to specify the same successive
centralizers.  Derived restriction is transitive, so the derived
fiber obtained from $M_{L'}$ is the same as that from $M_L$.
Induction proves the assertion for this orbit.

Consider an orbit for which no such torus exists.  The direction
$\lambda-h_e/2$ centralizes the triple, so it is central in $L$.
Since $\lambda$ is central in $H_L$, this implies that $H_L$
centralizes $h_e$.  In the centralizer of $h_e$, an element
centralizing $e$ also centralizes $f$, by uniqueness of the
compatible triple with $e,h_e$ fixed.  Thus $S_e$ equals the
triple centralizer in $H_L$. It is reductive and finite modulo
$Z(L)^\circ$. Since $\lambda-h_e/2$ is central in $L$, every vector
in $E_L$ has $h_e$-weight two. Lemma~\ref{amp:orbits} gives only
nonpositive weights in the normal space, so that space is zero:
$H_Le$ is open in $E_L$. The irreducible space $E_L$ has only one
open orbit.

Every boundary orbit has already been handled by the proper-Levi
case.  Apply Lemma~\ref{amp:support} to those finitely many orbits
with tests $U\otimes\eta$.
Every neutral cocharacter of a triple in $L$ annihilates
$\eta\in X^*(L)$.  Their boundary contributions therefore have
upper bounds independent of $\eta$.
Together with Equation~\eqref{amp:levi-bound}, the localization
triangle shows that the open orbit alone satisfies the same type
of uniform bound for each fixed $U$.

We explain why these invariant bounds detect the entire corrected
fiber.  The group $S_e$ is reductive, possibly disconnected, and
contains the central torus $Z(L)^\circ$ with finite quotient.
Restrictions of characters of $L$ form a finite-index sublattice
of $X^*(Z(L)^\circ)$.  Modulo tensoring by these restrictions,
there are finitely many irreducible rational representations of
$S_e$: there are finitely many central-character cosets, and for each
fixed central character the irreducibles are representations of
a finite-dimensional twisted group algebra of the finite group
$S_e/Z(L)^\circ$.  Choose finitely many fixed
$H_L$-representations $U_1,\ldots,U_r$ whose restrictions detect
the duals of representatives of these classes.
Such representations exist by the matrix-coefficient and
complete-reducibility argument used above.

Arbitrary rational representations of a reductive group in
characteristic zero are direct sums of finite-dimensional
irreducibles \cite[Remark~12.53\textup{(a)}, Theorem~22.42, and
Corollary~22.43]{MilneAlgebraicGroups}.
Invariants are exact, so every nonzero cohomology
class of the corrected fiber is detected by one of the tests
$U_j\otimes\eta$.  The correction on that test has only the
$h_e$-weights of $U_j$, since $h_e$ annihilates $\eta$.
More explicitly, an irreducible constituent of $H^j(W_e)$ of
$\rho_e$-weight $b$ has a nonzero invariant pairing with one of
these tests.  If the chosen test component has $h_e$-weight $a$,
that pairing has ordinary degree $j+b+a$: the total $\rho_e$-weight
is zero, and the test's numerical degree plus its $\rho_e$-weight
is its $h_e$-weight.  Hence
\[
 j+b\leq C'(U_j)-\min\operatorname{wt}_{h_e}(U_j).
\]
The maximum over the finite list gives the bound in
Equation~\eqref{amp:corrected-fiber}.  This argument includes rank zero, where $E_L=0$ and all
representation types are characters.  It starts the induction
and completes the proof.
\end{proof}

\subsection{The extreme amplitude estimate}

Choose a rational Weyl-invariant positive-definite metric and use
it to identify weight and cocharacter directions.  Conjugate
$\lambda$ into the closed dominant chamber.  For sufficiently
divisible positive integers $m$, the weight $\mu=m\lambda$ is
integral, dominant, and trivial on the connected center.  Let
$V_\mu$ be the corresponding irreducible representation.

\begin{theorem}[Extreme upper amplitude]
\label{amp:amplitude}
In the finite Koszul specialization at $[t]$ of
$\mathsf K(V_\mu^{\boxtimes d})$, retain the part with maximal
absolute value for the $d$-fold slide at the chosen leg.
Its cohomological upper bound is
\begin{equation}
\label{amp:amplitude-bound}
 C+\max_{e\in E_t}d\langle\mu,h_e\rangle.
\end{equation}
The constant is independent of $m$.
It is also uniform when $t$ is replaced by $tz$, with
$z\in Z(\Gd)^\circ$, after omitting the absolute-value central
part, and when one chooses among a fixed finite list of
hyperspecial variants.  Fixed additional kernels at disjoint
places are permitted.

Unless $t$, modulo its absolute-value central part, is a compact
factor times a Jacobson--Morozov point for $Q$, the difference
between $2d\langle\lambda,\mu\rangle$ and the nonconstant term
in Equation~\eqref{amp:amplitude-bound} is positive and grows
linearly with $m$.
\end{theorem}

\begin{proof}
In the loop description, the representation label becomes
$V_\mu^{\otimes d}$ and the $d$-fold slide acts by its holonomy;
this is the slide convention of Proposition~\ref{cat:loops}.
Among the weights of $V_\mu$, pairing with $\lambda$ is uniquely
maximal at $\mu$.  Indeed a weight lies in the convex hull of the
Weyl orbit of $\mu$, and the linear functional defined by $\mu$
attains its maximum on that orbit only at $\mu$ itself.
This argument also applies on chamber walls.
Thus the extreme part is the line of character $d\mu$.
It is $H$-stable; its character is trivial on the derived subgroup
of $H$, and the unipotent part of holonomy acts trivially on it.
Nilpotent thickenings in the Koszul specialization do not change
this semisimple eigenvalue separation.

Apply Proposition~\ref{amp:fiber} and
Lemma~\ref{amp:support} with this one-dimensional representation.
There are finitely many orbits by Lemma~\ref{amp:orbits}.
The maximum of their fixed costs is a single constant $C$,
and the representation contribution on the orbit of $e$ is
$d\langle\mu,h_e\rangle$.  This proves
Equation~\eqref{amp:amplitude-bound}.

Compatibility of the triple implies that
$\lambda-h_e/2$ centralizes it.  The invariant metric extends,
on the derived Lie algebra, to an invariant bilinear form;
therefore
\[
 (h_e,\lambda-h_e/2)
 =([e,f],\lambda-h_e/2)=0.
\]
Consequently
\begin{equation}
\label{amp:squared-gap}
 2d\langle\lambda,\mu\rangle
   -d\langle\mu,h_e\rangle
 =2dm\|\lambda-h_e/2\|^2.
\end{equation}
All directions here are projected away from the connected center.
If the right side vanishes, the semisimple element
$tQ^{-h_e/2}$ centralizes the triple and has absolute values one
in the remaining directions.  It is conjugate into a maximal
compact subgroup after removing the absolute-value central part
\cite[Theorems~3.6 and~3.7]{AdamsTaibiCohomology}.
This is exactly the required compact-times-Jacobson--Morozov form.
Conversely that form gives equality for its triple.
If equality occurs for no orbit, the finite list of squared norms
in Equation~\eqref{amp:squared-gap} has a strictly positive minimum.
This is the asserted uniform linear gap.

Finally, replacing $t$ by $tz$ leaves $H$, $E_t$, their orbits,
the allowed Levis, and every flag test unchanged.  It also leaves
the semisimple direction $\lambda$, the rational deshearing,
and all corrected-fiber and representation-test gradings unchanged,
even when $\log_Q|z|$ is irrational.  The only change
in the finite invariant Koszul test is in the scalar values
subtracted from a fixed list of invariant generators.
Its length, shifts, and the preceding geometric bounds are
unchanged.  All constants in the slice, flag, support, and
finite-representation arguments can therefore be chosen uniformly
on this central family.  The label $\mu$ is central-trivial, so
its extreme line is unchanged as well.
For finitely many hyperspecial variants take the maximum of their
constants.  Kernels at other places were fixed throughout
Proposition~\ref{amp:tests} and do not affect this uniformity.
\end{proof}
